\section{Discussion}

The results of the load and stress tests provide an empirical evaluation of the technical trade-offs between Edge and Cloud architectures. The following analysis validates the stability of the \textit{Thin Edge} design within the Plant Up! ecosystem.

\subsection{Validating the ``Thin Edge'' Architecture}

To evaluate the chosen architecture, the benefits of centralized cloud logic were compared with the response times of localized edge hardware. Results from Scenario 1 indicate that centralizing logic in the cloud is a viable strategy for Plant Up!. While the Edge-Centric control variant (Architecture B) processed thresholds locally in under 50\,ms, the Cloud-Centric variant (Architecture A) required over 200\,ms due to network latency and database processing. 

However, since botanical monitoring tracks slowly changing metrics, millisecond-level responsiveness is not a critical requirement. The cloud-based latency is generally not perceptible to the user. This justifies the use of centralized backend analytics and development agility over localized autonomy. Furthermore, centralizing decision-making within Supabase enables features such as the AI-driven Daily Mentor. Ultimately, the structural benefits of the Thin Edge approach outweigh the latency considerations for this application.

\subsection{Database Performance and Resilience}

The architecture must demonstrate the capacity to handle traffic spikes without service interruptions to ensure backend stability. In Scenario 2, the database architecture was evaluated using a load spike of 1,000 concurrent writes. The results indicated a peak ingestion throughput of approximately 850 writes/sec, with the internal queue normalizing within 1.2 seconds following the spike. 

This confirms that the Supabase environment, utilizing the TimescaleDB extension, maintains structural resilience during traffic bursts. Since Plant Up! does not require synchronous industrial control loops, the asynchronous latency associated with cloud-based evaluation is acceptable. The system maintained overall stability while absorbing the simulated load.

\subsection{Complexity vs. Operational Maintenance}

The architectural design also involves a balance between hardware autonomy and long-term maintainability. IoT literature suggests that centralizing logic can reduce developmental complexity, although it may increase dependency on network reliability \cite{zeetimThickThin, outsourceThinThick, deepwikiSupabase}. 

Maintaining conditional logic simultaneously in firmware (C++ for the ESP32) and in the cloud (SQL or Edge Functions) can introduce state synchronization challenges. The evaluation suggests that the Supabase backend is sufficiently resilient to support a Thin Edge approach. This allows Plant Up! to utilize simpler device firmware, where modifications to moisture thresholds or gamification algorithms require only a single backend deployment. This approach minimizes the risks associated with over-the-air (OTA) updates and supports rapid feature iteration.

\subsection{Architecture Limitations and Offline Resilience}

A fundamental limitation of employing a purist "Thin Edge" architecture is the strict dependency on continuous network availability. In the Plant Up! ecosystem, if the ESP32 microcontroller loses its local Wi-Fi connection, the firmware is intentionally designed to abort telemetry generation. It lacks a localized RAM buffer for caching metrics and instead immediately reverts to an Access Point (AP) setup mode. This visually signals to the user via LED indicators that the device requires network reconfiguration. 

Consequently, any environmental fluctuations occurring during this offline window are completely lost to the cloud backend. While data loss is generally discouraged in industrial IoT systems, it is an accepted and calculated trade-off within this context. Botanical parameters, such as soil moisture and light exposure, inherently change at a very slow pace. Missing a few hours of granular data has practically zero impact on overall plant health algorithms. Furthermore, the decision to prioritize standard Wi-Fi protocols over Bluetooth Low Energy (BLE) was precisely to guarantee direct device-to-cloud streams without localized gateway phones. Therefore, defaulting to a network-discovery mode gracefully handles outages without violating the system boundaries.

\subsection{Cost vs. Scalability Considerations}

Beyond technical throughput, evaluating an IoT architecture requires an analysis of its financial scalability. The current demonstration of Plant Up! operates entirely within the automated Free Tier of the Supabase managed platform. This provides a generous baseline for prototyping, rapid iteration, and academic evaluation without incurring direct monthly costs. 

However, as the hypothetical user base scales—for example, approaching 10,000 active devices—the data influx and database storage requirements would exceed these initial limits. At this medium scale, upgrading to the managed Supabase "Pro" subscription remains the most logistically sound strategy. The nominal cost of the premium tier easily outweighs the engineering labor hours required to provision and secure a custom database server. 

Moreover, Supabase circumvents the traditional pitfall of Backend-as-a-Service (BaaS) architectures: Vendor Lock-in. Because the entire platform is structured as a collection of open-source components explicitly wrapping a standard PostgreSQL database, the Plant Up! backend is highly portable. Should managed hosting costs ultimately become restrictive at a theoretical enterprise scale, the entire database alongside the TimescaleDB extensions can be seamlessly migrated. Utilizing Docker, it could be self-hosted on bare-metal infrastructure or dedicated Instances across AWS or Azure, proving the architecture is as economically scalable as it is technically capable.
